\chapter{总结、局限与后续工作}
\section{已完成工作的总结}
本工程已形成面向 1080p 视频的 HDMI 输入、帧自适应 Gamma、RGB565 DDR3 帧缓存、
逐帧姿态快照、定点镜头/姿态逆投影、双槽分块渲染和 HDMI 输出的完整 RTL 通路。
板载 CP2102 支持三个独立功能开关和亮度统计遥测，配置在输入帧边界原子生效。

当前 2026-10-22 基线的完整系统行为模型在给定并发与 DDR 停顿条件下，
达到约 30 个新结果帧/秒；最终 PDS 双角时序及 FIFO Gray 偏斜审计通过。
这些结论支持当前设计在 RTL 与实现工具层面的可行性。
\fillin{补齐实板测量后的最终性能、画质与演示结论}。

\section{赛题完成情况与当前局限}
\begin{table}[htbp]\centering\small
\caption{最终提交前的完成情况核对模板}
\begin{tabularx}{\textwidth}{p{3.1cm}YY}
\toprule
事项 & 当前边界 & 正式提交时补齐 \\
\midrule
视频完整链路 & RTL、实现报告已有；当前候选未板测 & \fillin{板测记录及演示视频} \\
标定参数 & 固定确认的内参与系数 & \fillin{原始标定图集及参数来源说明} \\
切向畸变 & 参数接口保留，当前值为零 & \fillin{非零场景测试或适用范围说明} \\
任意角度能力 & yaw/pitch 每轴约 $\pm15^\circ$，无 roll & \fillin{范围验收、安装标定与扩展状态} \\
重采样与黑边 & 最近邻分块、无效域置黑 & \fillin{双线性需求、画质与黑边验收} \\
实时性与延迟 & 行为模型约 30 fps，实板延迟未测 & \fillin{实板新帧率与端到端数据} \\
图像增强 & 七档帧自适应 Gamma & \fillin{实际光照实验} \\
资源、时序和 CDC & 当前版本有归档结果 & \fillin{最终提交版本与报告对应关系} \\
\bottomrule
\end{tabularx}
\end{table}

运行时串口开关并不是相机内参/畸变参数在线加载接口；当前标定参数由固定配置集成。
如现场要求重新标定相机或修改镜头，应说明参数导出、量化、重新构建及验证过程。

\section{可进一步改进的方向}
后续可围绕分块双线性采样、roll 参与的三轴姿态矩阵、图像/IMU 时间同步、
运行时标定参数加载和更高吞吐开展扩展。稀疏几何网格或倒数 LUT 可作为
精度与资源折中的研究方向，但它们尚不是当前默认路径的已验证功能。

每项改进应分别给出数学误差、边界行为、存储需求和完整系统性能，再重新进行
时序与实板验收。当前最直接的后续工作是完成候选烧录、真实图卡测量和长时稳定性记录。

\guide{正式结论保留已验证且与选题有关的贡献及数值；
删除已完成的待补项，仍未完成的功能如实写入适用范围，不得改写为已达到赛题全部要求。}
